\chapter{Evaluation, Benchmarks, and In-Context Learning}
\label{ch:evaluation}

% Part II, Chapter 10: Measuring model capability without fooling yourself.
%
% Core question: Your model's validation loss is 2.8. What does that mean?
% Almost nothing -- because loss doesn't directly map to any human-relevant
% capability. So how do you actually know if your model is good?
%
% This chapter connects back to Project 1's "metrics can lie" insight
% (Lab 5 Exp 0, Ablation 1) and extends it to the evaluation of large models.
%
% Planned sections:
%   10.1  Hook: "A Lower Loss Does Not Mean a Better Model."
%   10.2  Perplexity and Its Limitations
%   10.3  Benchmark Suites: MMLU, HellaSwag, HumanEval, GSM8K
%   10.4  Few-Shot Evaluation Protocol
%   10.5  In-Context Learning: What It Is and Why It Matters
%   10.6  Contamination Detection: When Benchmarks Lie
%   10.7  Generation-Based Evaluation
%   10.8  Human Evaluation and Preference
%   10.9  Failure Modes in Evaluation
%
% Key thesis: Evaluation is the most underrated and most unsolved problem
% in the LLM pipeline. Industry spends enormous engineering resources on
% eval infrastructure, yet there is still no consensus metric for "model
% quality." Students must learn to design eval thoughtfully, not just
% report numbers.
%
% In-context learning is covered here (not in scaling) because it is both
% an emergent capability and an evaluation protocol. Few-shot prompting
% is how we measure model capability, not just how we use models.
%
% Lab 10 (planned): Design and run a mini eval suite for a small pretrained
% model. Compare perplexity, few-shot accuracy, and generation quality.
% Show that rankings can change depending on the eval protocol.

\textit{This chapter is a placeholder. It will cover perplexity limitations,
benchmark suites, few-shot evaluation, in-context learning, contamination
detection, generation-based and human evaluation, and evaluation failure
modes.}
